\chapter{Type System and Formal Mathematical Foundation}
\label{ch:type_system}

\section{Introduction to the Formal Engine}
\label{sec:type_engine_intro}

In this chapter, we formulate the mathematical and type-theoretic foundations of \ToposLang{} and Geometric Native Coding. We construct the formal type system, specify its static typing judgements, define its small-step operational semantics, and establish rigorous mathematical proofs of its core safety and concurrency theorems.

The formalization proceeds in five sequential phases:
\begin{enumerate}[leftmargin=1.5cm]
    \item \textbf{Discrete Algebraic-Topological Foundations:} Graded cell complexes, chain groups, boundary homomorphisms, combinatorial homology, and the nilpotency axiom ($\partial^2 = 0$).
    \item \textbf{Continuous Clifford-Algebraic Foundations:} Construction of Clifford algebras $\mathcal{C}\ell(p,q,r)$, geometric products, versor sandwich operators, and projective/conformal geometric algebras.
    \item \textbf{The Integrated Synthesis:} Multivector chains, cellular sheaves, and the dual-tier discrete-continuous concurrency operator.
    \item \textbf{Combinatorial Hodge Theory:} Combinatorial Hodge Laplacians ($L_k$), discrete Hodge decomposition, and invariant Betti number computation.
    \item \textbf{Formal Type System and Operational Semantics:} Typing judgements, small-step transition relations, and mechanized-grade proofs of Progress, Preservation, Spatial Lock-Free Commutativity, and Homotopy Nilpotency Preservation.
\end{enumerate}

---

\section{Discrete Algebraic-Topological Foundations}
\label{sec:discrete_topology}

\subsection{Graded Cell Complexes}
\begin{definition}[Directed Acyclic Cell Complex]
A \textbf{Directed Acyclic Cell Complex (DACC)}, denoted by $\mathcal{K} = (K, \partial)$, is a topological space partitioned into a graded disjoint union of finite cell sets:
\begin{equation}
    K = \bigsqcup_{k=0}^n K_k = K_0 \sqcup K_1 \sqcup K_2 \sqcup \dots \sqcup K_n
\end{equation}
where $n \in \mathbb{N}$ denotes the maximal topological dimension of the complex. Each element $e_i^k \in K_k$ is an oriented open $k$-dimensional topological cell homeomorphic to the open $k$-ball $\mathbb{B}^k = \{ \mathbf{x} \in \mathbb{R}^k \mid \|\mathbf{x}\| < 1 \}$.
\end{definition}

\begin{definition}[Chain Groups and Chain Spaces]
For each dimension $k \in \{0, 1, \dots, n\}$, the \textbf{$k$-th chain group} $C_k(\mathcal{K}; \mathbb{F})$ is the finite-dimensional vector space of formal linear combinations generated by the $k$-cells of $K_k$ over a field $\mathbb{F}$:
\begin{equation}
    c = \sum_{i=1}^{|K_k|} a_i e_i^k \quad \text{where} \quad e_i^k \in K_k, \; a_i \in \mathbb{F}
\end{equation}
For exact structural topology in \ToposLang{}, the coefficient field is the field of rational numbers $\mathbb{F} = \mathbb{Q}$. When continuous physical fields are embedded, $\mathbb{F}$ is extended to a Clifford multivector space $\mathcal{C}\ell(p,q,r)$.
\end{definition}

\subsection{The Boundary Homomorphism and Incidence Matrices}
\begin{definition}[Boundary Operator]
The \textbf{boundary operator} is a graded family of linear homomorphisms:
\begin{equation}
    \partial_k : C_k(\mathcal{K}; \mathbb{F}) \longrightarrow C_{k-1}(\mathcal{K}; \mathbb{F}) \quad (k \ge 1)
\end{equation}
mapping each $k$-cell $e_j^k \in K_k$ to an oriented formal linear combination of its $(k-1)$-dimensional boundary faces:
\begin{equation}
    \partial_k(e_j^k) = \sum_{i=1}^{|K_{k-1}|} [e_i^{k-1} : e_j^k] \, e_i^{k-1}
\end{equation}
where $[e_i^{k-1} : e_j^k] \in \{-1, 0, +1\}$ denotes the topological incidence number reflecting the relative orientation between face $e_i^{k-1}$ and cell $e_j^k$. For $k=0$, we define $\partial_0 \equiv 0$.
\end{definition}

\begin{definition}[Boundary Incidence Matrix]
Choosing fixed ordered bases for $K_{k-1}$ and $K_k$, the boundary homomorphism $\partial_k$ is uniquely represented by a sparse integer incidence matrix $B_k \in \mathbb{Z}^{|K_{k-1}| \times |K_k|}$:
\begin{equation}
    (B_k)_{ij} = [e_i^{k-1} : e_j^k]
\end{equation}
The boundary of an arbitrary chain $c \in C_k$ with coefficient vector $\mathbf{a} \in \mathbb{F}^{|K_k|}$ is computed via standard matrix-vector multiplication:
\begin{equation}
    \partial_k(c) \longleftrightarrow B_k \mathbf{a}
\end{equation}
\end{definition}

\begin{axiom}[Fundamental Boundary Nilpotency]
\label{ax:nilpotency}
A Directed Acyclic Cell Complex $\mathcal{K}$ is topologically well-formed if and only if the composition of consecutive boundary homomorphisms evaluates identically to the zero map across all dimensions:
\begin{equation}
    \partial_{k-1} \circ \partial_k \equiv 0 \quad \forall k \ge 2
\end{equation}
In matrix representation, boundary nilpotency is equivalent to the matrix identity:
\begin{equation}
    B_{k-1} B_k = \mathbf{0}_{|K_{k-2}| \times |K_k|} \quad \forall k \in \{2, \dots, n\}
\end{equation}
\end{axiom}

\begin{remark}[Geometric and Systems Intuition: Why Boundaries Have No Boundaries]
Geometrically, Axiom \ref{ax:nilpotency} formalizes the universal physical property that \textit{a boundary has no boundary of its own}:
\begin{itemize}[leftmargin=1.2cm]
    \item The boundary of a solid 2D disk $\mathbb{D}^2$ is its 1D circular perimeter $S^1$; the boundary of a closed circle is empty ($\partial S^1 = \emptyset$).
    \item The boundary of a solid 3D ball $\mathbb{B}^3$ is its 2D spherical surface $S^2$; the boundary of a closed sphere is empty ($\partial S^2 = \emptyset$).
\end{itemize}
In systems programming, boundary nilpotency is the geometric equivalent of \textbf{pointer closure and structural integrity}. An operation that produces $\partial^2 \neq 0$ corresponds to creating an unclosed, torn geometric manifold with dangling references or orphaned boundary nodes. By validating $\partial^2 = 0$ statically at compile time, \ToposLang{} completely prevents structural state tearing before code generation begins.
\end{remark}

\subsection{Cycles, Boundaries, and Combinatorial Homology}
A sequence of chain groups and boundary operators satisfying Axiom \ref{ax:nilpotency} constitutes a chain complex:
\begin{equation}
    0 \longleftarrow C_0 \xleftarrow{\; \partial_1 \;} C_1 \xleftarrow{\; \partial_2 \;} C_2 \xleftarrow{\; \partial_3 \;} \dots \xleftarrow{\; \partial_n \;} C_n \longleftarrow 0
\end{equation}

\begin{definition}[Cycles and Boundaries]
For each dimension $k$:
\begin{itemize}[leftmargin=1.5cm]
    \item The subspace of \textbf{$k$-cycles}, denoted by $Z_k(\mathcal{K})$, is the kernel of $\partial_k$:
    \begin{equation}
        Z_k(\mathcal{K}) = \kerop(\partial_k) = \{ c \in C_k \mid \partial_k c = 0 \}
    \end{equation}
    \item The subspace of \textbf{$k$-boundaries}, denoted by $B_k^{\text{top}}(\mathcal{K})$, is the image of $\partial_{k+1}$:
    \begin{equation}
        B_k^{\text{top}}(\mathcal{K}) = \im(\partial_{k+1}) = \{ c \in C_k \mid \exists c' \in C_{k+1}, \; c = \partial_{k+1} c' \}
    \end{equation}
\end{itemize}
By Axiom \ref{ax:nilpotency}, every boundary is necessarily a cycle: $B_k^{\text{top}}(\mathcal{K}) \subseteq Z_k(\mathcal{K})$.
\end{definition}

\begin{definition}[Homology Groups and Betti Numbers]
The \textbf{$k$-th homology group} $H_k(\mathcal{K}; \mathbb{Q})$ is the quotient vector space:
\begin{equation}
    H_k(\mathcal{K}; \mathbb{Q}) = \frac{Z_k(\mathcal{K})}{B_k^{\text{top}}(\mathcal{K})} = \frac{\kerop(\partial_k)}{\im(\partial_{k+1})}
\end{equation}
The \textbf{$k$-th Betti number} $\beta_k$ is the dimension of the $k$-th homology group:
\begin{equation}
    \beta_k = \dim_{\mathbb{Q}}(H_k(\mathcal{K}; \mathbb{Q})) = \dim(\kerop(B_k)) - \rank(B_{k+1})
\end{equation}
Computationally:
\begin{itemize}[leftmargin=1.5cm]
    \item $\beta_0$ counts the number of connected components in the memory complex;
    \item $\beta_1$ counts independent non-contractible 1-cycles (circular information loops or deadlocks);
    \item $\beta_2$ counts enclosed 2-dimensional voids or trapped volumetric domains.
\end{itemize}
\end{definition}

\begin{definition}[Euler Characteristic]
The \textbf{Euler characteristic} $\chi(\mathcal{K})$ is the alternating sum of cell counts, which by the classical Euler-Poincaré theorem equals the alternating sum of Betti numbers:
\begin{equation}
    \chi(\mathcal{K}) = \sum_{k=0}^n (-1)^k |K_k| = \sum_{k=0}^n (-1)^k \beta_k
\end{equation}
\end{definition}

\subsection{Downward Support Closures}
To reason about the spatial memory footprint of computational operations, we define topological closures.

\begin{definition}[Downward Support Closure]
For an individual cell $e_j^k \in K_k$, its downward closure $\cl(e_j^k) \subset K$ is the set containing the cell itself and all of its iterated boundary faces down to dimension zero:
\begin{equation}
    \cl(e_j^k) = \{ e_j^k \} \cup \bigcup_{e_i^{k-1} \in \partial(e_j^k)} \cl(e_i^{k-1})
\end{equation}
For an arbitrary formal chain $c = \sum a_i e_i^k \in C_k$, its closure is the union over all basis cells with non-zero coefficients:
\begin{equation}
    \cl(c) = \bigcup_{\{i \mid a_i \neq 0\}} \cl(e_i^k)
\end{equation}
\end{definition}

\begin{definition}[Spatial Support of a Rewrite Rule]
Given a computational homotopy rewrite rule $r : c_{\text{lhs}} \Rightarrow c_{\text{rhs}}$, its total spatial support footprint $\supp(r) \subset K$ is the union of the downward closures of its input and output chains:
\begin{equation}
    \supp(r) = \cl(c_{\text{lhs}}) \cup \cl(c_{\text{rhs}})
\end{equation}
\end{definition}

---

\section{Continuous Clifford-Algebraic Foundations}
\label{sec:continuous_clifford}

\subsection{Construction of Clifford Algebras}
\begin{definition}[Clifford Algebra $\mathcal{C}\ell(p,q,r)$]
Let $V = \mathbb{R}^{p,q,r}$ be an $n$-dimensional real vector space ($n = p + q + r$) equipped with a symmetric bilinear form $\langle \cdot, \cdot \rangle$ having $p$ positive eigenvalues, $q$ negative eigenvalues, and $r$ zero eigenvalues. The \textbf{Clifford algebra} $\mathcal{C}\ell(p,q,r)$ is the associative unital algebra generated by $V$ subject to the fundamental Clifford relation:
\begin{equation}
    \mathbf{v}^2 = Q(\mathbf{v}) \mathbf{1} = \langle \mathbf{v}, \mathbf{v} \rangle \mathbf{1} \quad \forall \mathbf{v} \in V
\end{equation}
where $\mathbf{1}$ denotes the algebraic identity element.
\end{definition}

Let $\{e_1, e_2, \dots, e_n\}$ be an orthogonal basis of $V$ satisfying:
\begin{equation}
    e_i^2 = \begin{cases}
        +1 & \text{for } 1 \le i \le p \\
        -1 & \text{for } p+1 \le i \le p+q \\
        0  & \text{for } p+q+1 \le i \le p+q+r
    \end{cases}
    \quad \text{and} \quad
    e_i e_j = -e_j e_i \quad (i \neq j)
\end{equation}
The Clifford algebra $\mathcal{C}\ell(p,q,r)$ has dimension $2^n$ over $\mathbb{R}$. A general multivector $A \in \mathcal{C}\ell(p,q,r)$ is expressed as a linear combination of $2^n$ basis blades:
\begin{equation}
    A = \sum_{I \subseteq \{1,\dots,n\}} a_I e_I = a_0 \mathbf{1} + \sum_{i} a_i e_i + \sum_{i<j} a_{ij} e_{ij} + \dots + a_{1\dots n} e_{1\dots n}
\end{equation}
where $e_I = e_{i_1} e_{i_2} \dots e_{i_k}$ with $i_1 < i_2 < \dots < i_k$.

\subsection{The Geometric Product and Its Decomposition}
For any two vectors $\mathbf{u}, \mathbf{v} \in V$, the associative \textbf{geometric product} decomposes into the sum of a symmetric inner contraction and an antisymmetric outer wedge product:
\begin{equation}
    \mathbf{u}\mathbf{v} = \mathbf{u} \cdot \mathbf{v} + \mathbf{u} \wedge \mathbf{v}
\end{equation}
where:
\begin{align}
    \mathbf{u} \cdot \mathbf{v} &= \frac{1}{2}(\mathbf{u}\mathbf{v} + \mathbf{v}\mathbf{u}) = \langle \mathbf{u}, \mathbf{v} \rangle \in \mathbb{R} \quad \text{(grade 0 contraction)} \\
    \mathbf{u} \wedge \mathbf{v} &= \frac{1}{2}(\mathbf{u}\mathbf{v} - \mathbf{v}\mathbf{u}) \in \bigwedge^2 V \quad \text{(grade 2 bivector area)}
\end{align}

\subsection{Concrete Geometric Algebras in ToposLang}
\ToposLang{} natively integrates three specific Clifford signatures:

\subsubsection{3D Projective Geometric Algebra (PGA: $\mathbb{R}^{3,0,1}$)}
PGA introduces a single degenerate basis vector $e_0$ satisfying $e_0^2 = 0$. The multivector space has dimension $2^4 = 16$. PGA adopts a plane-based dual representation of Euclidean space:
\begin{itemize}[leftmargin=1.5cm]
    \item \textbf{Planes (Grade 1):} $P = n_x e_1 + n_y e_2 + n_z e_3 + d e_0$;
    \item \textbf{Lines (Grade 2):} $L = \mathbf{d} + \mathbf{m} e_0$, encoding Plücker line coordinates directly;
    \item \textbf{Points (Grade 3):} $X = x e_{023} + y e_{031} + z e_{012} + w e_{123}$.
\end{itemize}
Euclidean rigid body motions (rotations and translations) form the Lie group $\mathrm{SE}(3)$ and are represented by 8-dimensional even multivectors termed \textbf{motors} $M$:
\begin{equation}
    M = \exp\left(-\frac{1}{2}\theta B\right) = R + e_0 T
\end{equation}
Transformations are executed via sandwiching:
\begin{equation}
    A' = M A M^{-1} = M A \widetilde{M}
\end{equation}
where $\widetilde{M}$ is the multivector reversion of $M$.

\subsubsection{3D Conformal Geometric Algebra (CGA: $\mathbb{R}^{4,1,0}$)}
CGA introduces two basis vectors $e_+$ ($e_+^2 = +1$) and $e_-$ ($e_-^2 = -1$), mapped into two null basis elements:
\begin{equation}
    e_0 = \frac{1}{2}(e_- - e_+) \quad (\text{origin}), \qquad e_\infty = e_- + e_+ \quad (\text{point at infinity})
\end{equation}
satisfying $e_0^2 = 0$, $e_\infty^2 = 0$, and $e_0 \cdot e_\infty = -1$. CGA produces a 32-dimensional multivector space where spheres, planes, circles, and point-pairs are represented as single blades:
\begin{align}
    \text{Point: } P &= \mathbf{x} + \frac{1}{2}\|\mathbf{x}\|^2 e_\infty + e_0 \\
    \text{Sphere: } S &= P - \frac{1}{2}r^2 e_\infty
\end{align}
Intersections of geometric primitives evaluate in closed algebraic form via the outer wedge product:
\begin{equation}
    \text{Circle} = S_1 \wedge S_2
\end{equation}

---

\section{The Integrated Synthesis: Geometric Topos Computation}
\label{sec:integrated_synthesis}

\subsection{Multivector Chains}
We synthesize discrete algebraic topology and continuous Clifford geometry into a unified structure.

\begin{definition}[Multivector Chain]
A \textbf{multivector $k$-chain} $c \in C_k(\mathcal{K}; \mathcal{C}\ell(p,q,r))$ is a formal linear combination:
\begin{equation}
    c = \sum_{i=1}^{|K_k|} \mathbf{A}_i e_i^k \quad \text{where} \quad \mathbf{A}_i \in \mathcal{C}\ell(p,q,r), \; e_i^k \in K_k
\end{equation}
The boundary operator distributes linearly over multivector coefficients:
\begin{equation}
    \partial_k\left( \sum_{i=1}^{|K_k|} \mathbf{A}_i e_i^k \right) = \sum_{i=1}^{|K_k|} \mathbf{A}_i \partial_k(e_i^k)
\end{equation}
\end{definition}

\subsection{Cellular Sheaves}
\begin{definition}[Cellular Sheaf over a DACC]
A \textbf{cellular sheaf} $\mathcal{F}$ over a Directed Acyclic Cell Complex $\mathcal{K} = (K, \partial)$ is defined by:
\begin{enumerate}
    \item An association of a Clifford algebra stalk $\mathcal{F}(e) = \mathcal{C}\ell(p,q,r)$ to each cell $e \in K$;
    \item A family of restriction homomorphisms $\rho_{v \trianglelefteq e} : \mathcal{F}(e) \to \mathcal{F}(v)$ for each incidence $v \in \partial e$, satisfying the cocycle condition:
    \begin{equation}
        \rho_{u \trianglelefteq v} \circ \rho_{v \trianglelefteq e} = \rho_{u \trianglelefteq e}
    \end{equation}
\end{enumerate}
\end{definition}

\subsection{The Dual-Tier Concurrency Operator}
In classical systems, thread non-interference is evaluated purely at the discrete pointer level. In \ToposLang{}, non-interference is evaluated through a **dual-tier operator**:

\begin{definition}[Dual-Tier Interference Operator]
Let $r_1, r_2$ be two rewrite rules with discrete spatial supports $\supp(r_1), \supp(r_2) \subset K$ and continuous geometric multivector blade states $S(r_1), S(r_2) \in \mathcal{C}\ell(p,q,r)$. The \textbf{Interference Operator} is defined as:
\begin{equation}
    \Interf(r_1, r_2) = \Big( \supp(r_1) \cap \supp(r_2) \neq \emptyset \Big) \;\land\; \Big( S(r_1) \wedge S(r_2) \neq 0 \Big)
\end{equation}
Two operations are \textbf{spatially disjoint} if $\Interf(r_1, r_2) = \text{False}$.
\end{definition}

---

\section{Combinatorial Hodge Theory}
\label{sec:combinatorial_hodge}

To govern control flow via exact invariants, \ToposLang{} adapts discrete differential geometry over cell complexes \cite{lim2020hodge}.

\begin{definition}[Combinatorial Hodge Laplacian]
For each dimension $k \in \{0, 1, \dots, n\}$, the \textbf{combinatorial Hodge Laplacian} $L_k : C_k(\mathcal{K}; \mathbb{Q}) \to C_k(\mathcal{K}; \mathbb{Q})$ is defined as:
\begin{equation}
    L_k = \partial_k^* \partial_k + \partial_{k+1} \partial_{k+1}^*
\end{equation}
In matrix representation over the standard cell bases:
\begin{equation}
    L_k = B_k^T B_k + B_{k+1} B_{k+1}^T \in \mathbb{Q}^{|K_k| \times |K_k|}
\end{equation}
where $B_0 \equiv \mathbf{0}$ and $B_{n+1} \equiv \mathbf{0}$.
\end{definition}

\begin{theorem}[Discrete Hodge Decomposition Theorem]
\label{thm:hodge_decomp}
For every dimension $k$, the chain space $C_k(\mathcal{K}; \mathbb{Q})$ admits an orthogonal direct sum decomposition:
\begin{equation}
    C_k(\mathcal{K}; \mathbb{Q}) = \im(B_k^T) \oplus \im(B_{k+1}) \oplus \kerop(L_k)
\end{equation}
Furthermore, the kernel of the Hodge Laplacian is canonically isomorphic to the $k$-th homology group:
\begin{equation}
    \kerop(L_k) \cong H_k(\mathcal{K}; \mathbb{Q})
\end{equation}
Consequently, the $k$-th Betti number equals the nullity of $L_k$:
\begin{equation}
    \beta_k = \dim(\kerop(L_k)) = |K_k| - \rank(L_k)
\end{equation}
\end{theorem}

\begin{proof}
Let $c \in \kerop(L_k)$. Then:
\begin{equation}
    \langle c, L_k c \rangle = \langle c, B_k^T B_k c \rangle + \langle c, B_{k+1} B_{k+1}^T c \rangle = \| B_k c \|^2 + \| B_{k+1}^T c \|^2 = 0
\end{equation}
Since the inner product is positive-definite, this holds if and only if:
\begin{equation}
    B_k c = 0 \quad (\text{i.e., } c \in \kerop(B_k) = Z_k) \quad \text{and} \quad B_{k+1}^T c = 0 \quad (\text{i.e., } c \in \kerop(B_{k+1}^T) = \im(B_{k+1})^\perp)
\end{equation}
Therefore, $\kerop(L_k) = Z_k \cap (B_k^{\text{top}})^\perp$. Since $Z_k = B_k^{\text{top}} \oplus (Z_k \cap (B_k^{\text{top}})^\perp)$, projecting onto the quotient space yields the canonical vector space isomorphism:
\begin{equation}
    \kerop(L_k) \cong \frac{Z_k}{B_k^{\text{top}}} = H_k(\mathcal{K}; \mathbb{Q})
\end{equation}
which establishes the theorem.
\end{proof}

\begin{remark}[Systems Intuition of the Hodge Laplacian: Discrete Diffusion and Deadlock Detection]
In physical systems and distributed algorithms, the combinatorial Hodge Laplacian $L_k$ functions as a \textbf{discrete differential diffusion operator}:
\begin{itemize}[leftmargin=1.2cm]
    \item The downward component $B_k^T B_k$ evaluates the divergence of boundaries, measuring net localized flux entering or exiting $(k-1)$-dimensional faces.
    \item The upward component $B_{k+1} B_{k+1}^T$ evaluates the curl of coboundaries, measuring circulatory flow rotating around $(k+1)$-dimensional cells.
\end{itemize}
The harmonic kernel $\kerop(L_k)$ comprises computational chains that possess simultaneously \textbf{zero divergence and zero curl}. In fluid dynamics, harmonic 1-forms correspond to irrotational, incompressible fluid circulations around physical holes (e.g., flow around a cylinder or airfoil).

In concurrent and distributed computing, harmonic 1-chains represent \textbf{distributed information vortices and communication deadlocks}: closed cycles of tasks or agents where each waits on the next, circulating dependencies indefinitely without reaching a terminal sink. When \ToposLang{} evaluates an invariant loop $\texttt{until betti(K, dim=1) == 0}$, the runtime verifies that $\kerop(L_1) = \{\mathbf{0}\}$. When this harmonic space vanishes, the runtime possesses a mathematical guarantee that all circular dependencies have definitively contracted, resolving all system deadlocks without heuristic timeouts.
\end{remark}

---

\section{Formal Typing Rules and Operational Semantics}
\label{sec:type_system_formal}

\subsection{Type Syntax}
The syntax of types in \ToposLang{} is formally defined as:
\begin{equation}
    \tau ::= \texttt{cell}\langle k \rangle \mid \texttt{chain}\langle k, \mathbb{F} \rangle \mid \texttt{complex}\langle n \rangle \mid \texttt{motor}\langle p,q,r \rangle \mid \texttt{rewrite}\langle k \rangle \mid \texttt{unit}
\end{equation}
where $k, n \in \mathbb{N}$ denote dimensions, and $\mathbb{F} \in \{ \mathbb{Q}, \mathcal{C}\ell(p,q,r) \}$ denotes the coefficient algebra.

\subsection{Typing Judgements}
A typing environment $\Gamma$ maps identifiers to their types. We define the inference rules for \ToposLang{}:

\begin{gather}
\frac{e_j^k \in K_k}{\Gamma \vdash e_j^k : \texttt{cell}\langle k \rangle} \quad (\textsc{T-Cell}) \\[1.5ex]
\frac{\Gamma \vdash e_i^k : \texttt{cell}\langle k \rangle \quad \forall i \in \{1,\dots,m\} \quad a_i \in \mathbb{F}}{\Gamma \vdash \sum_{i=1}^m a_i e_i^k : \texttt{chain}\langle k, \mathbb{F} \rangle} \quad (\textsc{T-Chain}) \\[1.5ex]
\frac{\Gamma \vdash c : \texttt{chain}\langle k, \mathbb{F} \rangle \quad k \ge 1}{\Gamma \vdash \partial_k(c) : \texttt{chain}\langle k-1, \mathbb{F} \rangle} \quad (\textsc{T-Boundary}) \\[1.5ex]
\frac{\Gamma \vdash c_{\text{lhs}} : \texttt{chain}\langle k, \mathbb{F} \rangle \quad \Gamma \vdash c_{\text{rhs}} : \texttt{chain}\langle k, \mathbb{F} \rangle \quad \partial_k(c_{\text{lhs}}) = \partial_k(c_{\text{rhs}})}{\Gamma \vdash \texttt{rewrite}(c_{\text{lhs}} \Rightarrow c_{\text{rhs}}) : \texttt{rewrite}\langle k \rangle} \quad (\textsc{T-Rewrite}) \\[1.5ex]
\frac{\forall k \in \{2,\dots,n\}, \quad B_{k-1} B_k = \mathbf{0}}{\vdash \mathcal{K} : \texttt{complex}\langle n \rangle \text{ valid}} \quad (\textsc{T-ValidComplex})
\end{gather}

\subsection{Small-Step Operational Semantics}
Program execution is formalized as a transition between machine configurations:
\begin{equation}
    \langle \mathcal{K}, \sigma \rangle \longrightarrow \langle \mathcal{K}', \sigma' \rangle
\end{equation}
where $\mathcal{K} = (K, \partial)$ is the current DACC memory complex and $\sigma$ is the active instruction statement.

\begin{definition}[Cellular Homotopy Step Rule]
Let rule $r = \texttt{rewrite}(c_{\text{lhs}} \Rightarrow c_{\text{rhs}}) : \texttt{rewrite}\langle k \rangle$. The application of $r$ on complex $\mathcal{K}$ evaluates as:
\begin{equation}
    \frac{
        \begin{array}{c}
            \partial_k(c_{\text{lhs}}) = \partial_k(c_{\text{rhs}}) \quad 
            \alpha_r \notin K_{k+1} \\
            \partial_{k+1}(\alpha_r) \triangleq c_{\text{rhs}} - c_{\text{lhs}}
        \end{array}
    }{
        \langle \mathcal{K}, \texttt{apply } r \rangle \longrightarrow \langle \mathcal{K} \cup \{\alpha_r\}, \texttt{unit} \rangle
    } \quad (\textsc{Step-Homotopy})
\end{equation}
\end{definition}

\begin{definition}[Invariant Loop Step Rules]
The execution of invariant-driven loops is governed by two deterministic transition rules. 
Let $W \triangleq \texttt{until betti(dim=}k\texttt{) == } m \texttt{ \{ } S \texttt{ \}}$.
\begin{gather}
    \frac{\beta_k(\mathcal{K}) == m}{\langle \mathcal{K}, W \rangle \longrightarrow \langle \mathcal{K}, \texttt{unit} \rangle} \quad (\textsc{Step-LoopTerm}) \\[1.5ex]
    \frac{\beta_k(\mathcal{K}) \neq m}{\langle \mathcal{K}, W \rangle \longrightarrow \langle \mathcal{K}, S \,;\, W \rangle} \quad (\textsc{Step-LoopIter})
\end{gather}
\end{definition}

---

\section{Theorems and Proofs of Safety and Concurrency}
\label{sec:theorems_proofs}

We now prove the core mathematical properties guaranteeing memory safety and lock-free concurrency.

\subsection{Homotopy Nilpotency Preservation}
\begin{lemma}[Homotopy Nilpotency Preservation Lemma]
\label{lem:nilpotency_preservation}
Let $\mathcal{K} = (K, \partial)$ be a valid cell complex satisfying Axiom \ref{ax:nilpotency}. Let $r = \texttt{rewrite}(c_{\text{lhs}} \Rightarrow c_{\text{rhs}})$ be a well-typed rewrite rule satisfying $\partial_k(c_{\text{lhs}}) = \partial_k(c_{\text{rhs}})$. The updated cell complex $\mathcal{K}' = \mathcal{K} \cup \{\alpha_r\}$ produced by attaching cell $\alpha_r$ of dimension $k+1$ with boundary $\partial_{k+1}(\alpha_r) = c_{\text{rhs}} - c_{\text{lhs}}$ satisfies boundary nilpotency:
\begin{equation}
    \partial_k \circ \partial_{k+1} \equiv 0
\end{equation}
\end{lemma}

\begin{proof}
We evaluate the composition of boundary operators on the newly attached cell $\alpha_r$:
\begin{align}
    (\partial_k \circ \partial_{k+1})(\alpha_r) &= \partial_k\big( \partial_{k+1}(\alpha_r) \big) \\
    &= \partial_k\big( c_{\text{rhs}} - c_{\text{lhs}} \big)
\end{align}
By the linearity of the boundary homomorphism $\partial_k$:
\begin{equation}
    \partial_k(c_{\text{rhs}} - c_{\text{lhs}}) = \partial_k(c_{\text{rhs}}) - \partial_k(c_{\text{lhs}})
\end{equation}
By the typing premise of rule $r$ (\textsc{T-Rewrite}), $\partial_k(c_{\text{lhs}}) = \partial_k(c_{\text{rhs}})$. Therefore:
\begin{equation}
    \partial_k(c_{\text{rhs}}) - \partial_k(c_{\text{lhs}}) = \mathbf{0}_{C_{k-1}}
\end{equation}
For any cell $e \in K_{k+1}$ previously in $\mathcal{K}$, $(\partial_k \circ \partial_{k+1})(e) = 0$ holds by the inductive validity of $\mathcal{K}$. Furthermore, since no $(k+2)$-cells have been attached, $\partial_{k+1} \circ \partial_{k+2}$ remains identically zero. Thus, $\mathcal{K}'$ satisfies boundary nilpotency across all dimensions.
\end{proof}

\subsection{The Spatial Lock-Free Commutativity Theorem}
\begin{theorem}[Spatial Lock-Free Commutativity Theorem]
\label{thm:spatial_commutativity}
Let $\mathcal{K}$ be a valid cell complex. Let $r_1 : c_{\text{lhs},1} \Rightarrow c_{\text{rhs},1}$ and $r_2 : c_{\text{lhs},2} \Rightarrow c_{\text{rhs},2}$ be two well-typed rewrite rules. If the spatial support closures of $r_1$ and $r_2$ are disjoint:
\begin{equation}
    \supp(r_1) \cap \supp(r_2) = \emptyset
\end{equation}
then the execution of $r_1$ and $r_2$ commutes identically:
\begin{equation}
    r_1 \circ r_2 \equiv r_2 \circ r_1
\end{equation}
and both rules may be executed concurrently across physical threads without locks or synchronization primitives.
\end{theorem}

\begin{proof}
Let the memory state of complex $\mathcal{K}$ be decomposed into:
\begin{enumerate}
    \item A cell dictionary $\mathcal{D} : \text{CellID} \to \text{CellData}$;
    \item Sparse incidence matrices $B_k \in \mathbb{Q}^{|K_{k-1}| \times |K_k|}$.
\end{enumerate}
The set of memory locations read or modified by an operation $r$ is bounded by its spatial support closure:
\begin{equation}
    \text{MemoryFootprint}(r) \subseteq \bigcup_{e \in \supp(r)} \Big( \{ \mathcal{D}[e] \} \cup \{ B_k[i, j] \mid e \in \{e_i^{k-1}, e_j^k\} \} \Big)
\end{equation}
Assume $\supp(r_1) \cap \supp(r_2) = \emptyset$. It follows that:
\begin{equation}
    \text{MemoryFootprint}(r_1) \cap \text{MemoryFootprint}(r_2) = \emptyset
\end{equation}
The write-sets and read-sets of $r_1$ and $r_2$ are strictly disjoint subsets of machine memory. By Bernstein's Concurrency Conditions:
\begin{equation}
    (W_1 \cap R_2) \cup (R_1 \cap W_2) \cup (W_1 \cap W_2) = \emptyset
\end{equation}
Executing $r_1$ followed by $r_2$ produces the complex $\mathcal{K}_{12} = \mathcal{K} \cup \{\alpha_{r_1}, \alpha_{r_2}\}$. Executing $r_2$ followed by $r_1$ produces $\mathcal{K}_{21} = \mathcal{K} \cup \{\alpha_{r_2}, \alpha_{r_1}\}$. In the sparse incidence matrices, the column insertion for $\alpha_{r_1}$ and the column insertion for $\alpha_{r_2}$ modify independent column indices. Since matrix column updates commute under disjoint indices, the resulting matrix representations are identical:
\begin{equation}
    \mathcal{K}_{12} \equiv \mathcal{K}_{21}
\end{equation}
Hardware threads executing $r_1$ and $r_2$ in parallel access disjoint cache lines, producing zero cache-coherency invalidations and requiring zero synchronization locks.
\end{proof}

\subsection{Dual-Tier Concurrency Extension Theorem}
\begin{theorem}[Dual-Tier Concurrency Extension]
\label{thm:dualtier_commutativity}
Let $r_1$ and $r_2$ be two rewrite rules whose discrete supports share a set of boundary cells $C_{\text{shared}} = \supp(r_1) \cap \supp(r_2) \neq \emptyset$. If the continuous multivector blade states $S(r_1), S(r_2) \in \mathcal{C}\ell(p,q,r)$ satisfy:
\begin{equation}
    S(r_1) \wedge S(r_2) = 0
\end{equation}
and the updates to the shared boundary cells are additive linear combinations of the stalks, then $r_1$ and $r_2$ commute:
\begin{equation}
    r_1 \circ r_2 \equiv r_2 \circ r_1
\end{equation}
\end{theorem}

\begin{proof}
When $S(r_1) \wedge S(r_2) = 0$, the continuous physical domains embedded within the shared boundary cells do not intersect in the geometric algebra representation (the outer product measures the non-zero span of oriented subspaces; a zero wedge product indicates geometric linear dependence or non-intersection of orthogonal null elements). Because multivector addition in Clifford algebras is commutative ($A + B = B + A$), the simultaneous accumulation of continuous field values into the stalk $\mathcal{F}(c)$ for $c \in C_{\text{shared}}$ commutes identically. Hence, the operations execute concurrently via lock-free atomic SIMD accumulation.
\end{proof}

\subsection{Type Soundness: Progress and Preservation}
\begin{theorem}[Type Soundness]
\ToposLang{} satisfies the standard Wright-Felleisen type soundness properties \cite{wright1994syntactic}:
\begin{enumerate}
    \item \textbf{Preservation (Subject Reduction):} If $\vdash \mathcal{K} : \texttt{complex}\langle n \rangle$ is valid and $\langle \mathcal{K}, \sigma \rangle \longrightarrow \langle \mathcal{K}', \sigma' \rangle$, then $\vdash \mathcal{K}' : \texttt{complex}\langle n \rangle$ is valid.
    \item \textbf{Progress:} If $\vdash \mathcal{K} : \texttt{complex}\langle n \rangle$ is valid and $\sigma$ is a well-typed process statement, then either $\sigma$ is a value ($\texttt{unit}$) or there exists a configuration $\langle \mathcal{K}', \sigma' \rangle$ such that $\langle \mathcal{K}, \sigma \rangle \longrightarrow \langle \mathcal{K}', \sigma' \rangle$.
\end{enumerate}
\end{theorem}

\begin{proof}
\textbf{Preservation:} By induction on the derivation of $\langle \mathcal{K}, \sigma \rangle \longrightarrow \langle \mathcal{K}', \sigma' \rangle$. The only state-modifying rule is \textsc{Step-Homotopy}. By Lemma \ref{lem:nilpotency_preservation}, attaching cell $\alpha_r$ preserves boundary nilpotency ($\partial^2 = 0$). By Axiom \ref{ax:nilpotency} and rule \textsc{T-ValidComplex}, $\mathcal{K}'$ is a valid complex.

\textbf{Progress:} Let $\sigma$ be well-typed. If $\sigma = \texttt{apply } r$, then by rule \textsc{T-Rewrite}, $\partial_k(c_{\text{lhs}}) = \partial_k(c_{\text{rhs}})$. The premise of \textsc{Step-Homotopy} is satisfied, allowing a transition. If $\sigma$ is an invariant loop, the Betti numbers $\beta_k(\mathcal{K})$ are well-defined by Theorem \ref{thm:hodge_decomp} and evaluate deterministically via rational RREF. Either \textsc{Step-LoopTerm} or \textsc{Step-LoopIter} applies. Thus, well-typed configurations never become stuck.
\end{proof}

\subsection{The Euler--Poincar\'e Consistency Theorem}
\begin{theorem}[Euler--Poincar\'e Invariant Consistency]
\label{thm:euler_poincare}
For any well-typed Directed Acyclic Cell Complex $\mathcal{K} = (K, \partial)$, the alternating sum of cell counts identically equals the alternating sum of Betti numbers:
\begin{equation}
    \chi(\mathcal{K}) = \sum_{k=0}^n (-1)^k |K_k| = \sum_{k=0}^n (-1)^k \beta_k
\end{equation}
Furthermore, executing a cellular homotopy rewrite $r$ of dimension $k$ shifts the Euler characteristic by:
\begin{equation}
    \chi(\mathcal{K} \cup \{\alpha_r\}) = \chi(\mathcal{K}) + (-1)^{k+1}
\end{equation}
\end{theorem}

\begin{proof}
By the rank-nullity theorem applied to each linear boundary incidence matrix $B_k \in \mathbb{Q}^{|K_{k-1}| \times |K_k|}$:
\begin{equation}
    |K_k| = \dim(\kerop(B_k)) + \rank(B_k)
\end{equation}
By definition, $\beta_k = \dim(\kerop(B_k)) - \rank(B_{k+1})$. Rearranging terms yields:
\begin{equation}
    \dim(\kerop(B_k)) = \beta_k + \rank(B_{k+1})
\end{equation}
Substituting this expression into the cell count yields:
\begin{equation}
    |K_k| = \beta_k + \rank(B_{k+1}) + \rank(B_k)
\end{equation}
Now, evaluating the alternating sum of cell counts across all dimensions $k \in \{0, 1, \dots, n\}$ (with boundary conditions $\rank(B_0) = 0$ and $\rank(B_{n+1}) = 0$):
\begin{align}
    \sum_{k=0}^n (-1)^k |K_k| &= \sum_{k=0}^n (-1)^k \Big( \beta_k + \rank(B_{k+1}) + \rank(B_k) \Big) \\
    &= \sum_{k=0}^n (-1)^k \beta_k + \sum_{k=0}^n (-1)^k \rank(B_{k+1}) + \sum_{k=0}^n (-1)^k \rank(B_k)
\end{align}
Shifting the index of the second summation ($j = k+1$) produces:
\begin{equation}
    \sum_{k=0}^n (-1)^k \rank(B_{k+1}) = - \sum_{j=1}^{n+1} (-1)^j \rank(B_j) = - \sum_{k=0}^n (-1)^k \rank(B_k)
\end{equation}
The two rank summations cancel identically:
\begin{equation}
    \sum_{k=0}^n (-1)^k |K_k| = \sum_{k=0}^n (-1)^k \beta_k
\end{equation}
When a computational homotopy attaches cell $\alpha_r \in K_{k+1}$, the cell count $|K_{k+1}|$ increases by exactly 1, altering the Euler characteristic by $(-1)^{k+1} \cdot 1 = (-1)^{k+1}$. This establishes the invariant consistency of global topology across all homotopy execution steps.
\end{proof}

